%! Carrying Out a Test for the Difference of Two Population Proportions — Unit 6, Lesson 6.11 — Exit Ticket (Answer Key)
\documentclass[10pt]{article}
\usepackage{apstats-article}
\usepackage{apstats-key}   % pulls in -boxes; do NOT also load -boxes
\newcommand{\TallMath}[1]{$\displaystyle #1\rule[-1.4em]{0pt}{3.2em}$}

\begin{document}

\pageheader{Unit 6, Lesson 6.11}{Exit Ticket --- Answer Key}
\namedateperiod

\vspace{0.15in}

Juniors: 88/160, Seniors: 100/200. $\hat p_J=0.550$, $\hat p_S=0.500$, $\hat p_c\approx0.522$.
$H_a: p_J\ne p_S$. $\alpha=0.05$.

\textit{Do:}\\
$SE = \sqrt{0.522(0.478)(1/160+1/200)} = \sqrt{0.2495(0.01125)} = \sqrt{0.002807} \approx \ans{0.05298}$\\
$z = (0.550-0.500)/0.05298 = 0.050/0.05298 \approx \ans{0.94}$\\
$p\text{-value} = 2 \times P(Z > 0.94) \approx 2(0.1736) = \ans{0.347}$

\textit{Conclude:} \ansline{Since $p$-value $= 0.347 > \alpha = 0.05$, we fail to reject $H_0$. We do not have convincing evidence that the proportions of juniors and seniors who complete all assignments on time differ.}

\end{document}
